\chapter{Robotics Transformer Era}

\section*{학습 목표}

이 장은 Robotics Transformer 시대를 VLA의 탄생기로 정리한다. 핵심은 특정 논문 이름을 외우는 것이 아니라, 로봇 action이 transformer의 sequence prediction 문제로 재정의되었다는 점이다. 언어 문장을 token sequence로 보고 다음 token을 예측하듯이, 로봇 trajectory도 observation, language instruction, history가 주어졌을 때 다음 action 또는 action chunk를 예측하는 문제로 볼 수 있다. RT-1, PaLM-E, RT-2, RoboCat은 이 전환을 서로 다른 각도에서 밀어붙인 사례다 \cite{rt12022,palme2023,rt22023,robocat2023}.

\begin{enumerate}
  \item robot policy가 task-specific controller에서 generalist transformer policy로 이동한 배경을 설명한다.
  \item action tokenization과 multimodal sequence modeling의 의미를 정리한다.
  \item BC-Z, Gato, RT-1, PaLM-E, RT-2, RoboCat의 기술적 역할을 비교한다.
  \item 왜 이 흐름이 곧바로 완전한 Real VLA가 아니었는지, 다음 장의 action representation 문제로 연결한다.
\end{enumerate}

\section{전환점: action as sequence prediction}

고전적인 로봇 제어에서 action은 controller가 계산하는 입력이다. imitation learning에서는 expert trajectory에서 state-action pair를 모아 policy를 학습한다. Robotics Transformer 시대의 전환은 이 문제를 sequence modeling으로 재해석한 데 있다.
\begin{equation}
  p_{\theta}(\act_{1:T} \mid \obs_{1:T},\ell)
  =
  \prod_{t=1}^{T}
  p_{\theta}(\act_t \mid \obs_{\leq t},\act_{<t},\ell).
  \label{eq:robot-sequence-model}
\end{equation}
이 식은 behavior cloning의 일반화다. 차이는 입력과 출력이 더 이상 작은 vector pair가 아니라 multimodal token sequence가 된다는 점이다.

Transformer는 다음과 같은 token 집합을 attention으로 처리한다.
\begin{equation}
  X_t
  =
  [
  x^{text}_{1:N},
  x^{image}_{1:M},
  x^{state}_{1:R},
  x^{action}_{<t}
  ].
  \label{eq:robot-transformer-tokens}
\end{equation}
policy는 이 token sequence에서 action distribution을 예측한다.
\begin{equation}
  p_{\theta}(\act_t \mid X_t)
  =
  \mathrm{Head}_{act}
  \left(
  \mathrm{Transformer}_{\theta}(X_t)
  \right).
  \label{eq:robotics-transformer-action-head}
\end{equation}
이 구조는 language, vision, proprioception, action history를 하나의 sequence 처리 문제로 묶는다. VLA라는 이름은 여기서 자연스럽게 나온다. vision-language representation이 action head와 연결되기 때문이다.

\section{Language-conditioned imitation learning}

BC-Z류 흐름은 language-conditioned imitation learning이 generalization을 만들 수 있음을 보여주는 초기 단계로 볼 수 있다 \cite{bcz2022}. task id 대신 language instruction이나 video context를 condition으로 넣고, policy가 여러 task를 하나의 parameter set으로 학습한다.
\begin{equation}
  \mathcal{L}_{BC}
  =
  -
  \sum_{(\obs_t,\ell,\act_t^{E})\in\mathcal{D}}
  \log
  p_{\theta}(\act_t^{E}\mid \obs_t,\ell).
  \label{eq:language-bc-loss}
\end{equation}
이 식은 단순하지만 중요한 의미가 있다. task label이 language로 바뀌면, policy는 task 간 공유 구조를 language embedding을 통해 활용할 수 있다. ``open drawer'', ``close drawer'', ``pick red block'' 같은 명령이 서로 독립된 class가 아니라 text space 안에서 관련성을 갖는다.

그러나 이 단계의 한계도 분명하다. language-conditioned BC는 dataset coverage에 강하게 의존한다. unseen task generalization이 가능하더라도, 그 generalization은 observation distribution, object set, action interface, robot morphology가 크게 바뀌면 약해진다. 이 문제는 이후 RT-1과 RT-2에서 data scale과 model scale의 문제로 다시 등장한다.

\section{Gato와 generalist agent 관점}

Gato류 generalist agent 관점은 하나의 transformer가 서로 다른 modality와 task를 처리할 수 있다는 사고를 강화했다. 여기서 중요한 점은 robot-specific 성능 자체보다 interface 통일이다. text, image, discrete action, continuous value를 token처럼 정리하면 하나의 model family가 여러 domain을 다룰 수 있다.
\begin{equation}
  y_{k}
  \sim
  p_{\theta}(y_k \mid y_{<k}),
  \qquad
  y_k \in \mathcal{V}_{text}\cup\mathcal{V}_{image}\cup\mathcal{V}_{action}.
  \label{eq:generalist-token-model}
\end{equation}
이 관점은 로봇에 직접 적용할 때 곧바로 어려움에 부딪힌다. 로봇 action은 token prediction처럼 보이더라도 physical system에 작용한다. action error는 다음 token의 문법 오류가 아니라 collision, slip, instability로 나타난다. 그럼에도 generalist agent 관점은 VLA 연구에 중요한 질문을 남겼다. 하나의 model이 language understanding, visual recognition, robot control을 같은 sequence modeling framework 안에서 다룰 수 있는가?

\section{RT-1: real robot data와 action token}

RT-1은 Robotics Transformer라는 이름이 보여주듯, real robot manipulation data와 transformer policy를 결합한 대표 사례다 \cite{rt12022}. 핵심은 task diversity, large-scale robot trajectory, language instruction, action tokenization의 결합이다. policy는 observation과 instruction을 입력으로 받아 discretized action을 예측한다.
\begin{equation}
  \hat{a}_{t}^{disc}
  =
  \arg\max_{a\in\mathcal{A}_{disc}}
  p_{\theta}(a\mid \obs_{\leq t},\ell).
  \label{eq:rt1-discrete-action}
\end{equation}
discrete action token은 transformer training pipeline과 잘 맞는다. classification head를 붙이고 cross-entropy loss로 학습할 수 있기 때문이다.
\begin{equation}
  \mathcal{L}_{act}
  =
  -
  \sum_t
  \log
  p_{\theta}(a_t^{disc}\mid \obs_{\leq t},\ell).
  \label{eq:discrete-action-loss}
\end{equation}

이 구조는 VLA 관점에서 두 가지 의미가 있다. 첫째, language-conditioned robot policy가 large-scale data와 transformer architecture에서 실제 robot task를 수행할 수 있음을 보였다. 둘째, action representation이 model architecture만큼 중요하다는 사실을 드러냈다. action을 어떻게 discretize할지, time horizon을 어떻게 잡을지, control frequency를 어떻게 맞출지가 policy 성능과 안정성에 직접 영향을 준다.

\section{PaLM-E: embodied multimodal language model}

PaLM-E류 embodied multimodal language model은 language model 내부로 visual observation과 robot state를 넣는 방향을 보여준다 \cite{palme2023}. 이 접근은 robot policy를 단순한 action predictor가 아니라 embodied reasoning system으로 본다.
\begin{equation}
  h
  =
  \mathrm{LLM}_{\theta}
  (
  x^{text},
  E_{vision}(\obs_t),
  E_{state}(\conf_t)
  ).
  \label{eq:palm-e-embedding}
\end{equation}
여기서 $h$는 language response, plan, task-relevant representation, 또는 downstream action module의 입력이 될 수 있다.

이 흐름의 장점은 semantic reasoning이다. 대형 언어모델은 instruction decomposition, commonsense, long-horizon context 유지에 강하다. 하지만 action generation이 직접 연결되지 않으면 Chapter 10의 planner problem으로 돌아간다. 즉 embodied multimodal language model은 VLA의 reasoning backbone이 될 수 있지만, action interface와 control boundary를 별도로 설계해야 한다.

\section{RT-2: action을 token처럼 다루기}

RT-2류 접근은 VLM을 robot action과 결합한 전환점으로 볼 수 있다 \cite{rt22023}. 핵심 아이디어는 web-scale vision-language knowledge를 robot trajectory data와 함께 fine-tune하고, robot action을 text token과 유사한 방식으로 표현하는 것이다.
\begin{equation}
  p_{\theta}
  (y_{1:K}, a_{1:T}^{tok}
  \mid
  I,\ell)
  =
  \prod_k p_{\theta}(y_k \mid y_{<k},I,\ell)
  \prod_t p_{\theta}(a_t^{tok}\mid y_{1:K},a_{<t}^{tok},I,\ell).
  \label{eq:rt2-action-token}
\end{equation}
여기서 $y$는 language token, $a^{tok}$는 tokenized robot action이다. 이 구조의 의미는 크다. VLM이 갖고 있는 object, relation, semantic knowledge가 action prediction에 직접 연결될 가능성을 열었기 때문이다.

그러나 action-as-token에는 engineering 한계가 있다. discrete token은 training을 단순하게 만들지만, continuous control의 정밀도와 smoothness를 제한할 수 있다. 또한 token sequence가 길어지면 latency가 증가하고, high-frequency control에는 별도 controller가 필요하다. 따라서 RT-2 이후 연구는 action chunking, diffusion, flow matching, continuous action expert, efficient fine-tuning으로 확장된다. 이 내용은 Chapter 13에서 본격적으로 다룬다.

\section{RoboCat과 multi-embodiment adaptation}

RoboCat류 흐름은 multi-task, multi-embodiment, visual goal-conditioned policy, adaptation loop를 강조한다 \cite{robocat2023}. VLA에서 중요한 질문은 하나의 robot에서 많은 task를 하는 것만이 아니다. 다른 robot morphology, 다른 camera, 다른 object distribution으로 policy를 옮길 수 있는가가 중요하다.
\begin{equation}
  \pi_{\theta}
  :
  (\obs_t,\ell,e)
  \mapsto
  \act_t,
  \qquad
  e\in\mathcal{E}_{embodiment}.
  \label{eq:multi-embodiment-policy}
\end{equation}
여기서 $e$는 embodiment descriptor이다. robot이 바뀌면 action space, kinematics, camera geometry, controller dynamics가 바뀐다. 따라서 multi-embodiment policy는 단순한 multi-task policy보다 어렵다.

adaptation loop는 다음 형태로 볼 수 있다.
\begin{equation}
  \mathcal{D}_{k+1}
  =
  \mathcal{D}_{k}
  \cup
  \mathrm{Collect}
  (
  \pi_{\theta_k},\mathcal{T}_{new},\mathcal{E}_{new}
  ),
  \qquad
  \theta_{k+1}
  =
  \mathrm{Train}(\mathcal{D}_{k+1}).
  \label{eq:adaptation-loop}
\end{equation}
이 식은 data engine의 시작점이다. VLA는 model architecture만으로 성장하지 않는다. deployment, failure collection, retraining, evaluation loop가 함께 있어야 한다.

\section{Robotics Transformer 계열 비교}

\begin{table}[h]
  \centering
  \caption{Robotics Transformer 시대의 주요 흐름}
  \begin{tabularx}{\linewidth}{p{0.16\linewidth}X X}
    \toprule
    흐름 & 핵심 기여 & Real VLA 관점의 한계 \\
    \midrule
    BC-Z & language/video-conditioned imitation learning과 unseen task generalization & dataset coverage와 robot/domain shift에 민감하다 \\
    Gato & 하나의 transformer로 여러 modality/task를 다루는 generalist agent 관점 & physical execution과 safety boundary가 약하다 \\
    RT-1 & real robot data, task diversity, transformer action policy & action discretization과 data scale에 의존한다 \\
    PaLM-E & embodied multimodal language model과 state/image/text 결합 & action generation과 controller interface는 별도 설계가 필요하다 \\
    RT-2 & VLM knowledge와 robot action token을 결합한 VLA 전환점 & continuous/dexterous control, latency, action precision 문제가 남는다 \\
    RoboCat & multi-task, multi-embodiment adaptation과 data loop & embodiment transfer와 evaluation protocol이 어렵다 \\
    \bottomrule
  \end{tabularx}
  \label{tab:robotics-transformer-era}
\end{table}

이 표에서 보듯 Robotics Transformer 시대의 성과는 ``transformer를 로봇에 적용했다'' 정도로 요약하면 부족하다. 더 정확한 요약은 다음과 같다. robot policy가 language-conditioned, multimodal, scalable sequence model로 재정의되었다.

\section{왜 이것이 Real VLA의 시작인가}

VLA를 다음 mapping으로 쓰자.
\begin{equation}
  \act_t
  =
  \pi_{\theta}
  (\obs_t,\ell,h_t).
  \label{eq:vla-short-map}
\end{equation}
Robotics Transformer 시대는 이 mapping을 large-scale transformer로 학습할 수 있다는 가능성을 열었다. 하지만 Real VLA는 이 식 하나로 끝나지 않는다. 실제 시스템은 action representation, controller interface, data engine, evaluation, safety monitor를 필요로 한다.

따라서 이 시대는 첫 번째 완성형이면서 동시에 두 번째 문제 제기의 출발점이다. RT-1/RT-2 계열은 VLA를 가능하게 했지만, 다음 질문들을 남겼다.
\begin{enumerate}
  \item action을 discrete token으로 둘 것인가, continuous vector로 둘 것인가?
  \item 한 시점 action을 예측할 것인가, action chunk를 예측할 것인가?
  \item deterministic head를 쓸 것인가, distributional generator를 쓸 것인가?
  \item controller가 추적할 reference를 낼 것인가, hardware에 가까운 command를 낼 것인가?
  \item web-scale semantic prior가 physical skill로 전이되는 범위는 어디까지인가?
\end{enumerate}

\section{2024--2025 흐름으로 이어지는 질문}

RT-2 이후의 중요한 변화는 ``VLA를 만들 수 있다''에서 ``어떤 action interface와 fine-tuning recipe가 실제 robot에서 빠르고 안전한가''로 질문이 이동했다는 점이다. OpenVLA는 open-source 7B VLA와 대규모 real-world robot demonstrations를 통해 VLA 연구의 재현 가능한 기준점을 제공했다 \cite{openvla2024}. 그러나 open model 자체가 deployment-ready interface를 보장하지는 않는다.

OpenVLA-OFT는 같은 OpenVLA 계열에서도 action decoding, continuous action representation, action chunking, objective 선택이 fine-tuning 효율과 inference throughput에 큰 영향을 준다는 점을 보여준다 \cite{openvlaoft2025}. FAST는 action tokenization을 단순한 binning이 아니라 high-frequency robot action sequence를 압축하고 복원하는 문제로 재정의한다 \cite{fast2025}. $\pi_0$와 $\pi_{0.5}$는 pretrained VLM 위에 flow-based action generation과 heterogeneous co-training을 결합하여, action expert와 data mixture가 VLA generalization의 핵심 축임을 보여준다 \cite{pi02024,pi052025}.

\begin{warningbox}{이 장을 읽을 때의 critical reading point}
Robotics Transformer 계열을 읽을 때는 backbone 이름보다 action이 어떤 representation으로 기록되고, 어떤 controller interface로 실행되며, evaluation이 어떤 embodiment와 data split에서 이루어졌는지를 먼저 확인해야 한다. 최신 VLA 논문의 차이는 종종 architecture block 하나가 아니라 action tokenizer, chunk horizon, fine-tuning objective, data mixture, latency budget에 있다.
\end{warningbox}

\chapterwork
{BC-Z, Gato, RT-1, PaLM-E, RT-2, RoboCat을 history가 아니라 interface evolution으로 읽어라 \cite{bcz2022,gato2022,rt12022,palme2023,rt22023,robocat2023}.}
{Robotics Transformer 시대가 VLA를 가능하게 했지만 Real VLA를 완성하지는 못한 이유는 무엇인가?}
{RT-1, RT-2, OpenVLA 중 하나를 골라 input, action, controller, data, evaluation evidence를 Chapter 0 기준으로 분해하라.}

\section{요약}

Robotics Transformer 시대는 VLA의 역사에서 결정적이다. BC-Z는 language-conditioned imitation learning의 generalization 가능성을 보였고, Gato는 single transformer generalist agent 관점을 확산시켰다. RT-1은 real robot data와 transformer policy를 결합했고, PaLM-E는 embodied multimodal language model의 방향을 제시했다. RT-2는 VLM과 robot action token을 결합해 VLA라는 문제 설정을 선명하게 만들었고, RoboCat은 multi-task/multi-embodiment adaptation과 data loop를 강조했다.

하지만 이 장의 결론은 Robotics Transformer가 모든 문제를 해결했다는 것이 아니다. 오히려 이 시기 이후 가장 중요한 질문은 action representation으로 이동한다. action을 어떤 형식으로 표현하느냐가 control boundary, latency, dexterity, safety, data efficiency를 결정한다. 다음 장은 이 책의 기술적 중심부인 action representation으로 들어간다.
